% Source: https://artofproblemsolving.com/community/c3h2742533_confusing_problem
% Stable identity: aops-topic-2742533-q01
A number with at most $5$ divisors is $1$, a prime $p$, or of the form $p^2$, $p^3$, $pq$ or $p^4$ ($p,q$ distinct primes).

\emph{Small $N$ occur.} For $N=1,\dots,17$: $10,2,3,13,5,15,7,17,9,19,29,39,67,59,69,79,89$ have digit sums $1,\dots,17$, and each has at most $4$ divisors ($10=2\cdot5$, $15=3\cdot5$, $39=3\cdot13$, $69=3\cdot23$, $9=3^2$, the rest prime).

\emph{$N=18$ does not.} A number with digit sum $18$ is divisible by $9$. Among numbers with at most $5$ divisors, the multiples of $9$ are only $9=3^2$, $27=3^3$ and $81=3^4$ (any other multiple of $9$ has a further prime factor or a higher power of $3$, giving at least $6$ divisors). Their digit sums are all $9$, not $18$.

So $N=\boxed{18}$.
